\begin{abstract}
Los servicios de streaming musical generan enormes volúmenes de datos de
audio que ofrecen oportunidades para el descubrimiento automático de
patrones musicales. Este trabajo aplica algoritmos de agrupamiento
(\textit{clustering}) sobre el Spotify Tracks Dataset, que contiene
114,000 canciones de 125 géneros con 10 características de audio cuantitativas.
Se empleó K-Means como algoritmo principal, complementado con análisis PCA
para reducción dimensional y visualización. El proceso siguió la metodología
CRISP-DM. Los resultados identificaron [K] clusters con perfiles musicales
diferenciados, obteniendo un Silhouette Score de [valor], lo que evidencia
agrupaciones coherentes con los géneros musicales reales. La aplicación
práctica es un sistema de recomendación musical basado en similitud
de cluster.

\textbf{Palabras clave:} clustering, K-Means, Spotify, características de audio,
PCA, sistema de recomendación, CRISP-DM.
\end{abstract}
